\section{Identification of the driving process and the ordered curve}
\label{sec:identification}

We remain on the deterministic subsequence and fixed environment of
Lemma~\ref{lem:quenched-diagonal}.  Proposition~\ref{prop:moving-test}
now supplies the four-change crossing probability after localized
exploration.  To extract the Dobrushin driving process from this
probability, we insert a short wired interval on the free boundary.
The pairing martingale belongs to the resulting four-change law.
We first replace it by a bounded Loewner expression, and only then
compare that law with the original Dobrushin law.  This order is
what permits the interval to shrink without amplifying the
comparison error.  The method follows the \emph{Bounded martingale
tests} Lemma and \emph{Driving convergence} Proposition
of~\cite[Section ``Identification of the Dobrushin strand'']{C}.

\subsection{Bounded tests from a shrinking boundary interval}

Use the initial half-plane coordinate
$\psi_n=\Psi\circ\varphi_n$ from Section~\ref{sec:tightness},
reflecting the half-plane picture if necessary so that its positive
real side is primal-free.
Write $g_{n,t}$ and $W_n(t)$ for the maps and driver in
\eqref{eq:loewner-clock}.  Fix $T,R<\infty$ and stop at
$\tau_n=\tau_n(T,R)$.  Write $\E_n$ for expectation under the
original Dobrushin law in the fixed environment.  Let
$\mathcal F_{n,t}$ be the filtration of
the explored path and the switches it has encountered.  For
$x>2R+1$, define
\begin{equation}\label{eq:power-test}
 M_{n,t}(x)=
 \left(\frac{xg'_{n,t}(x)}{g_{n,t}(x)-W_n(t)}\right)^{1/2},
 \qquad 0\le t\le\tau_n.
\end{equation}
The real point $x$ is not swallowed on this range.  Indeed,
$g_{n,t}(x)\ge x$ and
\begin{equation}\label{eq:real-flow-bounds}
 \begin{split}
 x-R&\le g_{n,t}(x)-W_n(t)
       \le x+R+\frac{2T}{x-R},\\
 g'_{n,t}(x)&=
 \exp\left(-\int_0^t
       \frac{2\,du}{(g_{n,u}(x)-W_n(u))^2}\right).
 \end{split}
\end{equation}
Consequently $M_{n,t}(x)$ and its reciprocal are bounded by
constants depending only on $T,R,x$.

\begin{lemma}[Approximate martingale identities]
\label{lem:bounded-martingale-tests}
For each $x>2R+1$, every sequence of revealing stopping times
$0\le\sigma_n\le\upsilon_n\le\tau_n$, and every uniformly bounded
$\mathcal F_{n,\sigma_n}$-measurable sequence $F_n$, the original
Dobrushin law satisfies
\begin{equation}\label{eq:approximate-martingale}
 \E_n\left[
 F_n\bigl(M_{n,\upsilon_n}(x)-M_{n,\sigma_n}(x)\bigr)
 \right]\longrightarrow0.
\end{equation}
Stops during a known local arc are allowed.
\end{lemma}

\begin{proof}
Fix $s>0$, initially without taking a limit in $s$.  Insert a
separate primal wired interval on the formerly free side, with
lattice endpoints whose initial half-plane coordinates tend to
$x,x+s$.  In the open-cap four-change law, explore the strand
starting at the original initial mark $a_n$, whose half-plane
coordinate is $0$.  By
Lemma~\ref{lem:compatible-drawing}, the domain, initial charts,
and drawings may be chosen the same away from this insertion.
Regard the Dobrushin stopping rules and $F_n$ as functions of
revealed histories.  Apply the same rules on the common tree of
prefix histories, extend them as stopping rules outside that tree
by stopping at its first exit, and extend $F_n$ boundedly there.
These conventions retain the order of the two stops and their
localization; below we show that the relevant stopped prefixes
stay in the common tree.  Thus the same symbols will denote
well-defined stopping times and history tests under open caps.

We check the localization required by
Proposition~\ref{prop:moving-test}.  Real Loewner trajectories from
$x$ and $x+s$ remain separated from the driver, and their gap solves
\begin{equation}\label{eq:real-gap}
 \partial_t\bigl(g_{n,t}(x+s)-g_{n,t}(x)\bigr)
 =-\frac{2\bigl(g_{n,t}(x+s)-g_{n,t}(x)\bigr)}
 {(g_{n,t}(x+s)-W_n(t))(g_{n,t}(x)-W_n(t))}.
\end{equation}
For fixed $s$, this gap is bounded below by a positive constant.
The same ODE is regular in a complex neighborhood of the interval,
uniformly on the stopped range.  The prefix therefore stays out
of a fixed neighborhood of the insertion, and cannot terminate
at either inserted mark before the stop.  A point initially
corresponding to $iY$, with $Y$ sufficiently large, has a fixed
disk untouched by all these hulls.  It supplies the basepoint of
Proposition~\ref{prop:moving-test}.  The four marks, including tip
and infinity, are uniformly separated in the normalized charts
for each fixed $s$.  Initial marked Jordan convergence transfers
these statements to the physical drawings.  The history
tightness hypothesis follows from
Proposition~\ref{prop:curve-tightness} for the four fixed marks.

Let $Z_{n,t}^{s}$ be the conditional probability of the rare primal
pairing under the open-cap law.  It is an exact bounded
martingale in the revealed-switch filtration.  At each of the
allowed stopping times,
Proposition~\ref{prop:moving-test} gives
\begin{equation}\label{eq:four-mark-replacement}
 Z_{n,t}^{s}-f(\chi_{n,t}^{s})\longrightarrow0
 \quad\hbox{in probability and in mean},\qquad
 \chi_{n,t}^{s}=
 \frac{g_{n,t}(x+s)-g_{n,t}(x)}
      {g_{n,t}(x+s)-W_n(t)}.
\end{equation}
The vanishing initial errors in the lattice mark positions are
included here; they do not affect the fixed-$s$ limit.
Optional sampling and the boundedness of $F_n$ thus give the
approximate martingale identity for $f(\chi_{n,t}^{s})$.

Divide this identity by $f(s/(x+s))$ while $s$ is fixed.  The
small-interval behavior in
Proposition~\ref{prop:pure-four-mark}, together with analytic ODE
bounds near $x$, gives
\begin{equation}\label{eq:normalized-small-interval}
 \frac{f(\chi_{n,t}^{s})}{f(s/(x+s))}
       =M_{n,t}(x)+o_s(1)+o_n(1).
\end{equation}
Here $o_n(1)$ is taken first for fixed $s$, and $o_s(1)$ is uniform
over the stopped paths.  More explicitly,
\[
 \frac{\chi_{n,t}^{s}}{s/(x+s)}
 =\frac{xg'_{n,t}(x)}{g_{n,t}(x)-W_n(t)}+O(s),
\]
with the ratio bounded above and away from zero; applying
$f(u)\sim c\sqrt u$ proves the displayed replacement.  By
\eqref{eq:real-flow-bounds}, the expression after replacement is
bounded independently of small $s$.

We can now compare laws using only bounded functions of the
stopped histories.  First compare the deterministic separate
primal closure with the Dobrushin law.  The former adds just the
short wired interval, so FK monotonicity gives a coupling.
Explore its actual open clusters in the stronger law.  Unless
one reaches a fixed outer neighborhood of the insertion, the
configurations outside that neighborhood can be coupled
identically: the explored clusters are enclosed by closed cuts,
and the conditional outside laws then agree.  The insertion has
physical diameter tending to zero as $s\downarrow0$, by the
uniform continuity of the initial inverse maps.  The old primal
wire stays a positive distance away.  The intervening annulus
therefore has no primal wire attachments except at its rims,
and Lemma~\ref{lem:full-grid-arm} shows that the probability of
reaching the outer neighborhood tends to zero, first in the mesh
limit and then as $s\downarrow0$.

Second, compare this deterministic closure to open caps.
Lemma~\ref{lem:cap-change} gives a likelihood depending only on
the pairing, with two fixed positive values.  If the exceptional
pairing has probability $p$ under the separate-closure law,
normalizing these two weights
shows that the total variation difference is at most $C p$.
By Proposition~\ref{prop:moving-test} at the empty prefix and
Lemma~\ref{lem:cap-change}, for fixed $s$ that separate-closure
pairing probability tends to
$g(s/(x+s))$, which tends to zero as $s\downarrow0$.

The stopped prefixes avoid the spare neighborhood of the
insertion established above.  On the successful coupling event
their drawings, revealed histories, Loewner maps, and stopping
rules therefore agree.  The extension convention also makes
the stopped tests agree on these common histories.  Combining the two comparisons with
\eqref{eq:normalized-small-interval} yields
\[
 \limsup_{n\to\infty}
 \left|\E_n\!\left[
 F_n(M_{n,\upsilon_n}(x)-M_{n,\sigma_n}(x))\right]\right|
 \le e(s),\qquad \lim_{s\downarrow0}e(s)=0.
\]
This proves \eqref{eq:approximate-martingale}.  Notice that the
coupling errors were applied after replacement by bounded tests;
they were never divided by the small number $f(s/(x+s))$.
\end{proof}

\subsection{Two tests determine the driver}

\begin{proposition}[Driving convergence]
\label{prop:driving-convergence}
In the fixed environment under consideration, $W_n$ converges in
law, uniformly on compact capacity intervals, to
$\sqrt{16/3}\,B$, where $B$ is standard real Brownian motion.
\end{proposition}

\begin{proof}
We first prove tightness under the fixed localization $T,R$.
Set $h=1/2$ and, for $x>2R+1$, write
\[
 X_{n,t}(x)=g_{n,t}(x),\qquad
 A_{n,t}(x)=x^h g'_{n,t}(x)^h,\qquad
 M_{n,t}(x)=A_{n,t}(x)(X_{n,t}(x)-W_n(t))^{-h}.
\]
The real-flow bounds show that $X_n,A_n$ are uniformly bounded
and Lipschitz in time, with $A_n$ bounded away from zero.
Choose $2R+1<x_1<x_2$.  Equation~\eqref{eq:real-gap} also gives
a positive lower bound on
$X_{n,t}(x_2)-X_{n,t}(x_1)$ throughout the stopped range.

At a revealing stopping time $\sigma_n$, freeze these two
quantities and define, for $-R\le w\le R$,
\[
 F_j(w)=A_{n,\sigma_n}(x_j)
       (X_{n,\sigma_n}(x_j)-w)^{-h},\qquad j=1,2.
\]
Both derivatives are positive.  Their ratio is
\[
 r(w):=\frac{F'_1(w)}{F'_2(w)}
 =\frac{A_{n,\sigma_n}(x_1)}{A_{n,\sigma_n}(x_2)}
 \left(\frac{X_{n,\sigma_n}(x_2)-w}
            {X_{n,\sigma_n}(x_1)-w}\right)^{h+1}.
\]
The separation of the two real solutions and their uniform
bounds imply $r'(w)\ge c>0$.  Also $F'_2(w)$ is bounded above
and away from zero.  Put
$a_n=r(W_n(\sigma_n))$.  Integrating
$F'_2(w)(r(w)-a_n)$ on either side of $W_n(\sigma_n)$ gives
\begin{equation}\label{eq:two-test-convexity}
 F_1(w)-F_1(W_n(\sigma_n))
 -a_n\bigl(F_2(w)-F_2(W_n(\sigma_n))\bigr)
 \ge c_1|w-W_n(\sigma_n)|^2.
\end{equation}
The predictable coefficient $a_n$ is uniformly bounded.

Between $\sigma_n$ and another allowed stop $\upsilon_n$, changing
the frozen $X_n,A_n$ costs at most
$C(\upsilon_n-\sigma_n)$.  Apply
Lemma~\ref{lem:bounded-martingale-tests} to the two increments,
using the bounded $\mathcal F_{n,\sigma_n}$-measurable coefficient
$a_n$ as a test for the second.  Equation~\eqref{eq:two-test-convexity}
then yields
\begin{equation}\label{eq:driver-increment}
 \E_n|W_n(\upsilon_n)-W_n(\sigma_n)|^2
 \le C\E_n(\upsilon_n-\sigma_n)+o(1).
\end{equation}
This holds for every sequence of allowed stopping times.  In
particular, for the driver stopped at $\tau_n$, it implies the
stopping-time tightness criterion of Aldous~\cite{Aldous} in the
Skorokhod $J_1$ topology.
The drivers are bounded on this range.  Every prelimit driver is
continuous, and the continuous paths form a closed subspace for
$J_1$ convergence.  Thus all limits are continuous, and convergence
along a subsequence may be realized uniformly on compact stopped
intervals.

We identify any such limit inside an inner driver bound, choosing
inner levels which are continuity levels for the stopping operation.
These levels can be increased to the outer bound.  Continuity of
the Loewner ODE under uniform driver convergence gives limits
$X_t(x),A_t(x)$ and
\begin{equation}\label{eq:limiting-power-test}
 M_t(x)=A_t(x)(X_t(x)-W_t)^{-h}.
\end{equation}
Use bounded continuous functions of finitely many earlier driver
values as tests in \eqref{eq:approximate-martingale}, and then a
monotone-class argument.  It follows that each $M(x)$ is a
martingale in the limiting filtration.  This argument applies at
the localized stops as well.  The bounds above and below persist.

Before applying It\^o's formula, note that
\[
 W_t=X_t(x)-\left(\frac{A_t(x)}{M_t(x)}\right)^{1/h}.
\]
The processes $X,A$ are continuous of finite variation, $M$ is a
continuous martingale, and the function in this identity is smooth
on the range of their positive bounds.  Hence $W$ is a continuous
semimartingale.  Write its decomposition as a continuous local
martingale plus a continuous finite-variation process $B^{\rm fv}$,
and denote its quadratic variation by $[W]$.

Put $Z_t(x)=X_t(x)-W_t$.  The Loewner equations give
\[
 dX_t(x)=2Z_t(x)^{-1}\,dt,\qquad
 d\log A_t(x)=-2hZ_t(x)^{-2}\,dt.
\]
The finite-variation part of $dM_t(x)/M_t(x)$ is consequently
\[
 hZ_t(x)^{-1}\,dB_t^{\rm fv}
 -4hZ_t(x)^{-2}\,dt
 +\frac{h(h+1)}{2}Z_t(x)^{-2}\,d[W]_t.
\]
It vanishes because $M(x)$ is a martingale.  Multiplication by
$Z_t(x)^2/h$ gives an identity of signed measures,
\begin{equation}\label{eq:driver-characteristics}
 Z_t(x)\,dB_t^{\rm fv}-4\,dt
      +\frac{1+h}{2}\,d[W]_t=0.
\end{equation}
Subtract this identity for $x_1$ and $x_2$.  Their strictly
positive separation forces $dB^{\rm fv}=0$.  Substitution back
into \eqref{eq:driver-characteristics} yields
\[
 d[W]_t=\frac{8}{1+h}\,dt=\frac{16}{3}\,dt.
\]
L\'evy's characterization therefore identifies the stopped limit
as Brownian motion with variance rate $16/3$.  If the initial
half-plane picture was reflected, reflecting back negates the
driver and leaves this Brownian law unchanged.

It remains to remove the outer driver bound.  For any fixed
capacity horizon, the limiting upper bound on the probability
that $W_n$ exits $[-R,R]$ is bounded by the Brownian probability
of reaching inner levels tending to $R$.  This follows by applying
the stopped convergence at continuity levels below $R$; exiting
the outer interval entails first reaching each such inner level.
These Brownian probabilities tend to zero as $R\to\infty$.
Thus the full drivers are tight on every compact capacity
interval, and each subsequential law has the Brownian
identification just proved.  This proves the proposition.
\end{proof}

\subsection{From capacity time to the full oriented curve}

We can now complete the proof of Theorem~\ref{thm:main}.  We first
work in the common disk, and then return to the physical domains
and to environmental probability.

\begin{proof}[Proof of Theorem~\ref{thm:main}]
Fix the environment and subsequence of
Lemma~\ref{lem:quenched-diagonal}.  By
Proposition~\ref{prop:driving-convergence}, the drivers converge
uniformly on compact time intervals in law to
$W=\sqrt{16/3}\,B$.  For every fixed height $y>0$, continuity of
the inverse Loewner equation gives convergence of
\[
 H_n(t,y)=\Psi^{-1}\!\left(g_{n,t}^{-1}(W_n(t)+iy)\right)
\]
uniformly in $t$ on each compact capacity interval to the analogous
$H(t,y)$ for $W$.  A large driver bound has arbitrarily high
probability on that interval, so
Proposition~\ref{prop:tip-approximation} applies after localization.

For completeness, the uniform conclusion at height zero follows
without assuming uniform boundary convergence of the inverse
maps.  For two small positive heights, the triangle inequality
and \eqref{eq:tip-approximation} show that $H_n(\cdot,y)$ and
$H_n(\cdot,y')$ are close in the uniform metric, in probability,
after the mesh limit.  Passing at these fixed heights to the
limiting driver proves that $H(\cdot,y)$ is Cauchy in probability
in the space of continuous disk-valued functions on the time
interval as $y\downarrow0$.  Its limit is continuous.  The SLE
trace theorem~\cite{RS} identifies its pointwise values with the
SLE trace, first at rational times and then at all times by
continuity.  The same approximation argument now gives
\begin{equation}\label{eq:capacity-trace-convergence}
 \eta_n\ \Longrightarrow\ \eta
 \quad\hbox{uniformly on every compact capacity interval},
\end{equation}
where $\eta$ is chordal $\mathrm{SLE}_{16/3}$ in $\D$ from $-1$
to $1$.  In particular, a macroscopic excursion cannot be hidden
in a sequence of capacity intervals of vanishing length: the
uniform tip estimate already controls all times in those intervals.

To include the terminal part of the curve, first fix a small cap
about $1$.  Chordal SLE tends to its target almost surely
\cite{RS}; hence with probability as close to one as desired its
tail after a sufficiently large deterministic time lies in a
smaller cap.  By \eqref{eq:capacity-trace-convergence}, the
discrete curve has entered the original cap by that time with
the same high probability.  Lemma~\ref{lem:no-return} bounds
the diameter of its subsequent tail.  Letting the cap shrink
therefore makes both tails uniformly small.  Parametrize the
two curves by the common capacity clock until the chosen time
and then by any increasing parametrizations of their tails.
This proves convergence in the full uniform oriented-curve
metric modulo increasing reparametrization.

The initial inverse disk charts converge uniformly on their
closures by marked Jordan convergence.  Composing with them
transfers the curve convergence to the original physical
embedding.  Finally, the drawing comparison gives
$d(\gamma_n,\widehat\gamma_n)=O(\delta_n)$ in that same oriented
metric.  We have proved the asserted weak convergence of
quenched laws for every environment in the full-probability
event of Lemma~\ref{lem:quenched-diagonal}.

Return now to the entire prescribed mesh sequence.  The
oriented-curve space, with constant waiting intervals identified,
is separable, and its weak probability-law topology is metrized
by the bounded-Lipschitz distance for the bounded curve metric.
Each finite quenched law is measurable in its bonds.  From every
subsequence, Lemma~\ref{lem:quenched-diagonal} selects a further
deterministic subsequence along which this distance from the SLE
law tends to zero almost surely.  The subsequence criterion for
convergence in probability proves convergence in environmental
probability along the full sequence.  All choices requiring
small disorder were made before choosing the domain and its
approximation, as in the polygon comparison.  This completes
the proof.
\end{proof}
